\cotefile{157-2}
\scanpage{35}
\mspage{33}

le foncteur $(\sigma_M,\beta_M)$ précédent, i.e. qui donnent lieu à un \deleted{\unsure{diagramme}} \inserted{triangle} commutatif \unsure{(strict)}

\diagram{\begin{tikzcd}[ampersand replacement=\&]
 \& \underline{\mathrm{Ch}}_{\Sigma}(M) \arrow[d,"{(\sigma_M,\beta_M)}"] \\
 X \arrow[ur,"C"] \arrow[r,"{(F_0,F_1)}"] \& M\times M
\end{tikzcd}}

Mais si les foncteurs $F_0,F_1$ sont constants, de valeurs $a,b$ comme \unsure{tantôt|d'hab.}, i.e. $(F_0,F_1)$ constant de valeur $(a,b)$, alors \deleted{les} \inserted{la cat. des} foncteurs $C$ en question est \inserted{\deleted{\unsure{tau}} tautologiquement} isomorphe à celle des foncteurs de $\unsure{X}$ dans la fibre de $(\sigma_M,\beta_M)$ en $(a,b)$, d'où l'isomorphisme de catégories fondamental

\[
\boxed{\underline{\mathrm{Ch}}_{\Sigma_X}(X;a_X,b_X)\simeq
 \underline{\mathrm{Hom}}^{\unsure{!}}(X,\underline{\mathrm{Ch}}_{\Sigma}(M;a,b))}
\]

On va expliciter le contenu essentiel de cette formule en un

\uline{Scholie} \unsure{Soient} $M$ une catégorie, \unsure{avec|ou|soi} un ens $\Sigma$ de flèches de localisation, $a,b\in\mathrm{Ob}\,M$. Alors on peut trouver un \deleted{Ind} objet $\inserted{\underline{\underline K}=}K(a,b)$, \inserted{$K=(K_n)$} de $\mathrm{Ind}\,\mathrm{Cat}$\deleted{\unsure{, et}} (indexé par $\mathbb N$ mais peu importe au \unsure{fond|fait}) et un isom can, pour toute catégorie $X$ (isom fonctoriel en $X$)

\[
\begin{aligned}
\mathrm{Hom}_{\Sigma_X}(a_X,b_X)&\simeq
 \unsure{\Pi_0|\pi_0}\bigl(\underbrace{\underline{\mathrm{Hom}}(X,\underline{\underline K})}_{\substack{\Vert\,\unsure{\text{déf}}\\ \varinjlim\,\underline{\mathrm{Hom}}(X,K_n)}}\bigr)\\
 &\simeq\varinjlim\,\Pi_0\bigl(\underline{\mathrm{Hom}}(X,K_n)\bigr)\unsure{.}
\end{aligned}
\]

% archive sheet number: 34
